% Figure 2 — Ce que l'outil voit : chaque couche, et ce qu'elle couvre (ou pas),
% selon l'endroit où la commande agit.
\begin{tikzpicture}[
    cell/.style={draw=white,line width=1.2pt,minimum width=38mm,minimum height=12.5mm,
                 text width=35mm,align=center,font=\sffamily\scriptsize,inner sep=2pt},
    oui/.style={cell,fill=wepwnclair},
    part/.style={cell,fill=dASK!18},
    non/.style={cell,fill=grisclair,text=gris,
                postaction={pattern=north east lines,pattern color=trait}},
    ligne/.style={minimum width=34mm,minimum height=12.5mm,text width=32mm,
                  align=left,font=\sffamily\small,inner sep=2pt},
    tete/.style={minimum width=38mm,font=\sffamily\small\bfseries,align=center,text width=36mm},
  ]
  \def\W{3.85}\def\H{1.3}
  % En-têtes de colonnes
  \node[tete] at (1*\W,0.2) {Dans le projet};
  \node[tete] at (2*\W,0.2) {Sur la machine,\\hors du projet};
  \node[tete] at (3*\W,0.2) {Surfaces distantes\\{\scriptsize\mdseries base, cloud, git, SaaS}};

  % Lignes : (y, titre, sous-titre)
  \foreach \y/\t/\s in {1/Lire avant/le moteur (statique),
                        2/Récupérer/{instantané, point de contrôle},
                        3/Surveiller/{filet : ptrace ou WinFsp},
                        4/Refuser/plan de contrôle,
                        5/Demander/la personne (ask),
                        6/Tracer/les reçus}
    \node[ligne] at (0.15,-\y*\H) {\textbf{\t}\\{\scriptsize\color{gris}\s}};

  % Lire
  \node[oui]  at (1*\W,-1*\H) {prédit ce qui est détruit,\\créé, déplacé};
  \node[oui]  at (2*\W,-1*\H) {prédit la destruction\\$\rightarrow$ ESCALATE};
  \node[oui]  at (3*\W,-1*\H) {reconnaît la surface\\et la sorte d'effet};
  % Récupérer
  \node[oui]  at (1*\W,-2*\H) {copie exacte avant,\\tout ou rien};
  \node[non]  at (2*\W,-2*\H) {aucune copie\\(V17 : jamais prétendue)};
  \node[non]  at (3*\W,-2*\H) {aucune copie locale ;\\indice chez le fournisseur};
  % Surveiller
  \node[oui]  at (1*\W,-3*\H) {sauvegarde chaque fichier\\au moment où il est détruit};
  \node[part] at (2*\W,-3*\H) {enregistré, non sauvegardé\\(V18, filet Linux)};
  \node[non]  at (3*\W,-3*\H) {invisible au filet\\local};
  % Refuser
  \node[oui]  at (1*\W,-4*\H) {\texttt{.demo\_cli/}, \texttt{.demo\_cli.toml},\\\textit{hooks} du projet};
  \node[part] at (2*\W,-4*\H) {chemins listés seulement :\\\texttt{\textasciitilde/.bashrc}, \textit{hooks}, tâche};
  \node[non]  at (3*\W,-4*\H) {sans objet};
  % Demander
  \node[part] at (1*\W,-5*\H) {seulement si rien\\n'a pu être capturé};
  \node[part] at (2*\W,-5*\H) {destruction hors projet\\: ask};
  \node[oui]  at (3*\W,-5*\H) {ask ou deny\\selon la surface};
  % Tracer
  \node[oui]  at (1*\W,-6*\H) {chaque décision};
  \node[oui]  at (2*\W,-6*\H) {chaque décision};
  \node[oui]  at (3*\W,-6*\H) {chaque décision\\(sans secrets)};

  % Légende
  \node[oui,minimum width=6mm,text width=4mm,minimum height=4mm] (l1) at (0.4,-8.75) {};
  \node[right=1mm of l1,font=\sffamily\scriptsize] (t1) {couvert};
  \node[part,minimum width=6mm,text width=4mm,minimum height=4mm,right=5mm of t1] (l2) {};
  \node[right=1mm of l2,font=\sffamily\scriptsize] (t2) {partiel};
  \node[non,minimum width=6mm,text width=4mm,minimum height=4mm,right=5mm of t2] (l3) {};
  \node[right=1mm of l3,font=\sffamily\scriptsize] {non couvert --- dit tel quel dans la raison de la décision};
\end{tikzpicture}
